\begin{textsample}{POS Dim 3 – ac – Score 46.00 – ac/ac\_inf\_16.txt}  \label{ex:f3_pos_024}
16. \textbf{BEBÊS} ( ZERO A 1 ANO E 6 \textbf{MESES} ) 16.1.
Campo de experiências “ Espaços, tempos, quantidades, relações e transformações ” As \textbf{crianças} vivem inseridas em espaços e tempos de diferentes dimensões, em um mundo constituído de fenômenos naturais e socioculturais.
Desde muito \textbf{pequenas}, elas procuram se situar em diversos espaços ( rua, bairro, cidade etc. ) e tempos ( dia e noite ; hoje, ontem e amanhã etc. ). \textbf{Demonstram} também \textbf{curiosidade} sobre o mundo físico ( seu próprio corpo, os fenômenos atmosféricos, os animais, as plantas, as transformações da natureza, os diferentes tipos de materiais e as possibilidades de sua manipulação etc. ) e o mundo sociocultural ( as relações de parentesco e sociais entre as pessoas que conhece ; como vivem e em que trabalham essas pessoas ; quais suas tradições e seus costumes ; a diversidade entre elas etc. ).
Além disso, nessas experiências e em muitas outras, as \textbf{crianças} também se deparam, \textbf{frequentemente}, com conhecimentos matemáticos ( contagem, ordenação, relações entre quantidades, dimensões, medidas, comparação de \textbf{pesos} e de comprimentos, avaliação de distâncias, reconhecimento de formas geométricas, conhecimento e reconhecimento de numerais cardinais e ordinais etc. ) que igualmente aguçam a \textbf{curiosidade}.
Portanto, a Educação \textbf{Infantil} precisa promover experiências nas quais as \textbf{crianças} possam fazer observações, \textbf{manipular} objetos, investigar e explorar seu entorno, levantar hipóteses e consultar fontes de informação para buscar respostas às suas \textbf{curiosidades} e indagações.
Assim, a instituição escolar está criando oportunidades para que as \textbf{crianças} ampliem seus conhecimentos do mundo físico e sociocultural e possam utilizá -los em seu cotidiano. 16.2.
Direitos de aprendizagem e desenvolvimento - Conviver com seus pares e com \textbf{adultos}, explorando objetos e materiais de diferentes propriedades físicas, fazendo observações e construindo hipóteses sobre o mundo natural e social. - \textbf{Brincar} com seus pares e \textbf{adultos}, utilizando objetos, acessórios e elementos da natureza, assumindo diferentes papéis sociais e culturais, ampliando sua experiência relacional e \textbf{sensorial} com \textbf{texturas}, cheiros, cores, tamanhos, \textbf{pesos}, densidades e possibilidades de transformação. - Participar de situações de \textbf{cuidados} com o meio ambiente e sua sustentabilidade, de investigação sobre fenômenos naturais e resolução de problemas cotidianos que envolvam quantidades, medidas, dimensões, tempos, espaços, comparações, transformações, construindo hipóteses e explorando objetos e ferramentas diversas, tais como : bússola, lanterna, lupa, máquina \textbf{fotográfica}, smartphone, filmadora, gravador, projetor, computador e outras. - Explorar características do mundo natural e social, despertando a \textbf{curiosidade} por nomear, agrupar e ordenar informações para compreender os espaços à sua \textbf{volta} e as transformações neles realizados, sua própria história, os modos de vida das pessoas de sua comunidade e de outras culturas. - Expressar por meio das diferentes linguagens, suas impressões, observações, hipóteses e explicações sobre acontecimentos sociais, fenômenos da natureza e características do ambiente em que vive. - Conhecer -se e construir sua identidade pessoal e cultural, identificando seus próprios interesses na relação com o mundo físico e social, convivendo e conhecendo os costumes, as crenças e as tradições locais, regionais e de outras culturas. 16.2.1.
Objetivo de aprendizagem e desenvolvimento : ( EI01ET01 ) - Explorar e descobrir as propriedades de objetos e materiais ( odor, cor, sabor, temperatura ).
Aprendizagens específicas - Desenvolver a autonomia e o interesse por explorar e fazer descobertas. - Desenvolver o prazer da descoberta e da \textbf{curiosidade} em relação às propriedades de objetos e materiais, por meio de diferentes linguagens. - Construir, \textbf{progressivamente}, postura investigativa. - Apropriar -se, gradativamente, de conhecimentos sobre as características físicas e \textbf{sensoriais}, as propriedades e a utilidade dos objetos e materiais. - Construir hipóteses sobre as propriedades dos objetos e materiais.
Sugestões de experiências - Explorar e \textbf{manipular} objetos e \textbf{brinquedos} com diferentes \textbf{texturas}, cores, cheiros, \textbf{peso}, forma, espessura e outras características. - \textbf{Brincar} com elementos naturais ( água, \textbf{areia}, argila, folhas, pedras, sementes, etc. ). - Explorar diferentes objetos através do tato, investigando seus diferentes atributos : quente/frio, liso/áspero, grosso/fino, dentre outros. - Experimentar sabores e consistência ( sólido, pastoso e líquido ) de diversos alimentos ( picolé, carne, limão, mingau, suco, gelatina, frutas ). - \textbf{Brincar} de faz de conta, usando os objetos de forma convencional ou dando um novo significado por meio de sua \textbf{brincadeira} exploratória. 16.2.2.
Objetivo de aprendizagem e desenvolvimento : ( EI01ET02 ) - Explorar relações de causa e efeito ( transbordar, tingir, misturar, mover e remover etc. ) na interação com o mundo físico.
Aprendizagens específicas - Construir conhecimentos sobre os fenômenos do mundo físico. - Desenvolver a \textbf{curiosidade}, a concentração e a atenção por meio da exploração e interação com o mundo físico. - Desenvolver crescente autonomia e interesse pela exploração e investigação. - Perceber os processos de transformações simples. - Elaborar estratégias para compartilhar com seus pares o resultado de suas explorações. - Expressar seu prazer nas descobertas por meio das diferentes linguagens.
Sugestões de experiências - \textbf{Brincar} de encher e esvaziar um recipiente, \textbf{transferir} de um recipiente para outros de tamanhos diversos, utilizando elementos como : \textbf{areia}, água, pedregulhos e outros. - Experimentar misturas de \textbf{tintas} ( tinta guache e \textbf{tinta} produzida com elementos naturais ). - Realizar \textbf{pinturas}, utilizando \textbf{tintas} e misturas com : água, terra, argila, borra de café, óleo e outras possibilidades. - Participar de experiências com misturas diversas, utilizando água, terra, argila, óleo, sal, açúcar, borra de café, dentre outros, observando mudanças físicas e químicas. - \textbf{Brincar} com objetos, fazendo sombra ou luz com uso da lanterna. - Explorar objetos \textbf{empilhando}, segurando, jogando, retirando e guardando em uma caixa. - Observar alguns fenômenos da natureza como : chuva, vento, arco-íris, relâmpago, trovão, pôr do sol, dentre outros. - Ter contato com alguns fenômenos da natureza, por exemplo : colocar a mão para sentir os pingos de chuva, sentir o vento no rosto, tomar banho de sol e sentir seu calor. 16.2.3.
Objetivo de aprendizagem e desenvolvimento : ( EI01ET03 ) - Explorar o ambiente pela ação e observação, \textbf{manipulando}, experimentando e fazendo descobertas.
Aprendizagens específicas - Desenvolver, \textbf{progressivamente}, a postura investigativa. - Apropriar -se e construir \textbf{criativamente} significados sobre si e sobre o meio ambiente. - Construir, \textbf{progressivamente}, sentimento de pertença do meio ambiente. - Construir conhecimentos sobre plantas e animais no ambiente e seus modos de vida. - Construir, \textbf{progressivamente}, atitudes em relação a sustentabilidade, respeito e preservação ao meio ambiente. - Construir, gradativamente, conhecimentos científicos sobre os fenômenos físicos, químicos e biológicos nas relações com as experiências do cotidiano.
Sugestões de experiências - Explorar objetos e \textbf{brinquedos}, observando seu movimento e seus efeitos ( cata vento, queda de uma bola, bolinha de sabão, soprar bolinhas de isopor ou papel picado, puxar um \textbf{carrinho} com um barbante, etc. ). - Explorar o ambiente por meio dos sentidos : observar o barulho do vento, canto dos pássaros, cheiro das flores, \textbf{textura} das folhas, dentre outros. - \textbf{Manipular} materiais diversos ou \textbf{brinquedos}, provocando reações físicas como : movimento, inércia, flutuação, \textbf{equilíbrio} e outras. - \textbf{Brincar} com água, \textbf{areia} e terra usando tubos, baldes, peneiras, pá, rastelo, canecas e garrafas. - Participar de \textbf{brincadeiras} ao ar livre. - Deitar, se arrastar ou engatinhar na grama. - Participar da organização de cenários e ambientes estruturados. - Explorar o ambiente, por meio dos sentidos, observando e percebendo as características dos seres vivos. - Participar de práticas que estimulem a \textbf{curiosidade} sobre os elementos da natureza. 16.2.4.
Objetivo de aprendizagem e desenvolvimento : ( EI01ET04 ) - \textbf{Manipular}, experimentar, arrumar e explorar o espaço por meio de experiências de deslocamentos de si e dos objetos.
Aprendizagens específicas - Construir noção de espaço utilizando seu próprio corpo e objetos. - Expressar prazer em exercitar uma nova conquista motora, ao deslocar -se no espaço. - Construir, \textbf{progressivamente}, hipóteses quanto a coordenação de \textbf{gestos} e percepções ao realizar deslocamentos de si e dos objetos. - Desenvolver a orientação e percepção espacial, por meio das possibilidades de movimento e da exploração dos objetos no espaço. - Construir \textbf{hábitos} de \textbf{cuidados} e organização do espaço, materiais e objetos.
Sugestões de experiências - Deslocar -se no espaço, tentando alcançar objetos e \textbf{brinquedos} de seu interesse. - Participar de \textbf{brincadeiras} orientadas, de deslocamento de si e dos objetos ou \textbf{brinquedos}. - \textbf{Brincar} em diferentes posições : deitado, em cima, em baixo, de lado e outras possibilidades. - Explorar diferentes objetos : pegar, largar, chutar, \textbf{empilhar}, arremessar em várias direções e de diferentes modos, abrir e fechar, dentre outros. - Participar de \textbf{brincadeiras}, realizando deslocamentos, passando por obstáculos ( pneus, bambolês, caixas, cordas, cadeiras, mesas ). - Acompanhar com os olhos o movimento dos materiais e usar o corpo para explorar o espaço, \textbf{virando} para os diferentes lados ou rastejando -se. 16.2.5.
Objetivo de aprendizagem e desenvolvimento : ( EI01ET05 ) - \textbf{Manipular} materiais diversos e variados para comparar as diferenças e semelhanças entre eles.
Aprendizagens específicas - Desenvolver o interesse pela investigação dos diferentes atributos, por meio da manipulação e exploração de materiais diversos. - Perceber as diferenças e semelhanças entre os materiais. - Desenvolver diferentes estratégias para comparar, classificar, prever e descobrir características dos objetos e materiais diversos.
Sugestões de experiências - Coletar objetos diversos ( pedrinhas, \textbf{tampas} de garrafa, folhas, pedacinhos de madeira, papéis e outros ) para, com o apoio do \textbf{adulto}, organizar coleções, comparando as diferenças e semelhanças entre eles. - Atuar sobre materiais diversos, pegando, arremessando, rolando, apertando, cheirando, colocando um dentro do outro, em cima ou \textbf{batendo} um objeto no outro. - Experimentar gostos, sabores, \textbf{texturas}, odores e \textbf{sons}, de materiais diversos, realizando comparações simples entre eles. - \textbf{Brincar} com objetos e materiais variados, que produzam \textbf{sons}, refletem, ampliam, iluminam, dentre outras possibilidades. - \textbf{Brincar} com jogos de \textbf{encaixe}, \textbf{montando}, desmontando e construindo algo de seu interesse. - Explorar as características dos materiais fazendo uso das suas mãos, \textbf{pés}, \textbf{boca}, nariz e ouvido. - Participar da organização do espaço, guardando materiais semelhantes em uma caixa. 16.2.6.
Objetivo de aprendizagem e desenvolvimento : ( EI01ET06 ) - Vivenciar diferentes ritmos, velocidades e fluxos nas interações e \textbf{brincadeiras} ( em danças, balanços, escorregadores etc. ).
Aprendizagens específicas - Desenvolver, \textbf{progressivamente}, suas possibilidades corporais. - Ampliar \textbf{gestos}, postura e domínio do corpo, nas interações e \textbf{brincadeiras}. - Construir conhecimentos sobre a linguagem corporal ao explorar movimentos leves ou fortes, rápidos ou lentos, diretos, flexíveis. - Desenvolver a sociabilidade e a capacidade de se relacionar com o outro, construindo relações éticas, de respeito, tolerância, cooperação, solidariedade e confiança. - Apropriar -se \textbf{criativamente} da cultura de seu grupo, interagindo com seus pares, por meio de \textbf{brincadeiras}.
Sugestões de experiências - Acompanhar o ritmo da música marcando as batidas com \textbf{palmas} e os \textbf{pés}, aumentando ou diminuindo a velocidade. - \textbf{Brincar} no escorregador, gangorra e balanço, experimentando diferentes velocidades do seu deslocamento. - Explorar por meio da dança, movimentos leves ou fortes, rápidos ou lentos, sozinho ou interagindo com parceiros. - \textbf{Brincar} de se \textbf{balançar}, com apoio do \textbf{adulto}, em pneus amarrados por corda em troncos frondosos de árvore. - \textbf{Brincar}, com seus pares e \textbf{adultos}, de se \textbf{esconder} na área externa. - Participar de \textbf{brincadeiras} que envolvam o canto e o movimento, explorando \textbf{livremente} seu corpo. - Participar de \textbf{brincadeiras} que envolvam o canto e o movimento, ajustando seus \textbf{gestos} ao ritmo e versos da canção. 16.3.
Avaliação - \textbf{Acompanhamento} do processo de aprendizagem e desenvolvimento dos \textbf{bebês} - Observação \textbf{criteriosa}, \textbf{registro} e análise quanto : - Ao interesse por explorar e investigar as características físicas, propriedades e utilidades dos materiais e objetos do ambiente. - Às reações e iniciativas frente às transformações dos elementos físicos, químicos e naturais. - Às estratégias que elaboram para fazer comparações e identificar semelhanças e diferenças entre materiais diversos e variados. - Às estratégias que utiliza para explorar e orientar -se nos espaços. - Às iniciativas corporais frente aos diferentes ritmos, velocidades e fluxos nas interações e \textbf{brincadeiras}.

% matched lemmas: acompanhamento, adulto, areia, balançar, bater, bebê, boca, brincadeira, brincar, brinquedo, carro, criança, criativamente, criterioso, cuidado, curiosidade, demonstrar, empilhar, encaixe, equilíbrio, esconder, fotográfico, frequentemente, gesto, hábito, infantil, livremente, manipular, montar, mês, palma, pequeno, peso, pintura, progressivamente, pé, registro, sensorial, som, tampo, textura, tinta, transferir, virar, volta
\end{textsample}
